% reality/measurement_and_imperfect_data-DE.tex
\chapter{Messung und unvollkommene Daten}
\label{chap:measurement-and-imperfect-data}
\index{Beobachtung!Messdaten}
\index{Datenqualität!unvollkommene Daten}

Messung liefert Beobachtungsevidenz über realisierte Infrastruktur. Sie ersetzt
weder konstruktive Identität noch die Modellierung realisierter Zustände.

\section{Beobachtungskontext}
Messdaten in Eisenbahnprojekten sind hinsichtlich Auflösung, Genauigkeit,
Provenienz und Koordinatenkontext heterogen.
\begin{definition}[Beobachtete Daten]
\index{Beobachtete Daten}
Beobachtete Daten bezeichnen Messgrößen, die in einem bestimmten
Beobachtungskontext von Sensoren, Vermessungen oder historischen Datensätzen
gewonnen wurden.
\end{definition}
\begin{principle}[Prinzip des Beobachtungskontexts]
\index{Prinzip des Beobachtungskontexts}
Beobachtete Daten müssen gemeinsam mit Sensor-, Stationierungs-, Koordinaten-
und Topologiekontext interpretiert werden.
\end{principle}

\section{Eigenschaften unvollkommener Daten}
Beobachtete Daten enthalten typischerweise zufälliges Rauschen, systematische
Verzerrung, Unvollständigkeit und quellübergreifende Inkonsistenz.
\begin{definition}[Systematischer Fehler]
\index{Fehler!systematisch}
Ein systematischer Fehler ist eine dauerhafte Beobachtungsabweichung, die
durch Sensor-, Kalibrierungs- oder Transformationseinflüsse verursacht wird.
\end{definition}
\begin{definition}[Messunsicherheit]
\index{Unsicherheit!Messung}
Messunsicherheit quantifiziert das Intervall oder die Verteilung plausibler
physischer Werte, die mit einer Beobachtung vereinbar sind.
\end{definition}

\section{Inferenz aus Beobachtung}
Geometrie und Zustand werden durch Rekonstruktionsoperatoren und
Interpretationsregeln aus Beobachtungen erschlossen.
\begin{definition}[Diskrete Krümmungsschätzung]
\index{Krümmung!diskrete Schätzung}
Aus lokalen Richtungsstichproben kann eine diskrete Krümmungsschätzung
geschrieben werden als
\[
\kappa_i \approx \frac{\theta_{i+1}-\theta_i}{\Delta s_i}.
\]
\end{definition}
Solche Schätzungen hängen von der Abtastung ab und sind
unsicherheitsempfindlich; in der Praxis sind deshalb Filterung und
Regularisierung erforderlich.

\section{Beobachtung und Laufzeitrekonstruktion}
Beobachtung tritt über Laufzeitrekonstruktions-Pipelines in AIM ein, die Daten
mit metrischer Realisierung, Stationierungs- und Topologiekontext verbinden.
Beobachtete Werte werden als Evidenz für realisierten Zustand und nicht als
kanonische Objektidentität interpretiert.

\section{Ingenieurtechnische Verwendung}
In Engineering-Abläufen beeinflussen unvollkommene Daten Validierung,
Instandhaltungsplanung und Ziel--Realisiert-Vergleich. Robuste Abläufe
erfordern Provenienzverfolgung, unsicherheitsbewusste Bewertung und
ausdrückliche Verwaltung von Annahmen.
\index{Provenienz!Verfolgung}

\section{Verbindung zu AIM}
\begin{summary}
In AIM bilden Messung und unvollkommene Daten die Beobachtungsebene des Teils
Realität. Diese Ebene unterstützt realisierungsbewusste Rekonstruktion und
Interpretation und wahrt zugleich die Unterscheidung zwischen beobachteter
Evidenz und Engineering Identity.
\end{summary}
